\appendix{Представление графического материала}

Графический материал, выполненный на отдельных листах,
изображен на рисунках А.1--А.\arabic{числоПлакатов}.
\setcounter{числоПлакатов}{0}

\renewcommand{\thefigure}{А.\arabic{figure}} % шаблон номера для плакатов

\begin{landscape}
	
	\begin{плакат}
		
\includegraphics[width=0.82\linewidth]{pos1.eps}
		\заголовок{Сведения о ВКРБ}
		\label{pl1:image}
	\end{плакат}
	
	\begin{плакат}
		
\includegraphics[width=0.82\linewidth]{pos2.eps}
		\заголовок{Цель и задачи разработки}
		\label{pl2:image}
	\end{плакат}
	
	\begin{плакат}
		\includegraphics[width=0.82\linewidth]{pos3.eps}
		\заголовок{Диаграмма прецедентов для неавторизованного пользователя}
		\label{pl3:image}
	\end{плакат}
	
	\begin{плакат}
		\includegraphics[width=0.82\linewidth]{pos4.eps}
		\заголовок{Диаграмма прецедентов для заказчика}
		\label{pl4:image}
	\end{плакат}
	
	\begin{плакат}
		\includegraphics[width=0.82\linewidth]{pos5.eps}
		\заголовок{Диаграмма прецедентов для исполнителя}
		\label{pl5:image}
	\end{плакат}
	
	\begin{плакат}
		\includegraphics[width=0.82\linewidth]{pos6.eps}
		\заголовок{Уровневая схема архитектуры web-сервиса}
		\label{pl6:image}
	\end{плакат}
	
	\begin{плакат}
		\includegraphics[width=0.82\linewidth]{pos7.eps}
		\заголовок{Диаграмма последовательности обработки запроса}
		\label{pl7:image}
	\end{плакат}
	
	\begin{плакат}
		\includegraphics[width=0.82\linewidth]{pos8.eps}
		\заголовок{Реляционная модель базы данных}
		\label{pl8:image}
	\end{плакат}
	
	\begin{плакат}
		\includegraphics[width=0.82\linewidth]{pos9.eps}
		\заголовок{Заключение}
		\label{pl9:image}
	\end{плакат}
	
\end{landscape}